\documentclass{cssheet}

%--------------------------------------------------------------------------------------------------------------
% Basic meta data
%--------------------------------------------------------------------------------------------------------------

\title{Mathematische Aktivitäten in außerschulischen Lernorten}
\author{Prof. Dr. Christian Spannagel}
\date{\today}
\hypersetup{%
    pdfauthor={\theauthor},%
    pdftitle={\thetitle},%
    pdfsubject={Thema Masterarbeit},%
    pdfkeywords={master, phhd}
}

%--------------------------------------------------------------------------------------------------------------
% document
%--------------------------------------------------------------------------------------------------------------

\begin{document}

\vspace*{5mm}
\begin{center}
{\Large Thema für eine Masterarbeit}
\end{center}

\printtitle
\vspace*{1cm}

Außerschulische Lernorte bieten zusätzliche Zugänge zu mathematische Themen außerhalb des Klassenzimmers, die oftmals durch enaktives Material Hands-On-Erfahrungen möglichen. Durch diese Aktivitäten soll ein vertieftes Verständnis mathematischer Zusammenhänge erreicht und zur Beschäftigung mit mathematischen Sachverhalten motiviert werden.

In Zusammenarbeit mit der MAINS\footnote{Mathematik-Informatik-Station: \url{https://www.heidelberg-mains.org/}} sollen im Rahmen der Masterarbeit Materialien oder Workshops zu mathematischen Aktivitäten konzipiert, umgesetzt und evaluiert werden. Die Aktivitäten beziehen sich dabei auf ein spezielles Thema, z.\,B.
\begin{itemize}
\item Mathematik und Musik
\item Parkettierungen 
\item Platonische und Archimedische Körper
\item Spiegelungen
\item \ldots
\end{itemize}

Mögliche Tätigkeiten sind:
\begin{itemize}
\item \textbf{Literaturrecherche:} Aufarbeitung fachwissenschaftlicher und fachdidaktischer Literatur zu dem gewählte Thema, außerdem Recherche zu möglichen Hands-On-Aktivitäten
\item  \textbf{Umsetzung:} Produktion und Einsatz der Lernmaterialien oder Durchführung von Workshops mit Kindern und Jugendlichen im Kontext der MAINS 
\item  \textbf{Evaluation:} Beforschung der Wirksamkeit der Lernaktivitäten im Kontext einer passenden Forschungsfrage.
\end{itemize}

\vspace*{5mm}

\printlicense

\printsocials
\vspace*{5mm}



\end{document}
